%% Especificación de Requisitos de Software (SRS) — plantilla ISO/IEC/IEEE 29148
%% Compilar: tectonic -X compile srs.tex
\documentclass[11pt,a4paper]{article}
\usepackage{../latex/legasoft}

\lgtitulo{Especificación de Requisitos de Software}
\lgsubtitulo{Estructura conforme a ISO/IEC/IEEE 29148}
\lgtituloCorto{SRS}
\lgcodigo{LGS-REQ-001}
\lgversion{0.1}
\lgestado{Plantilla}
\lgfecha{\campo{fecha}}
\lgautor{\lgDirector\ (\lgDirectorRol)}

% --- Ficha reutilizable para un requisito funcional -------------------------
% Uso: \reqfun{RF-001}{Nombre}{Alta}{Descripción}{Criterio de aceptación}{Origen}
\newcommand{\reqfun}[6]{%
\par\vspace{4pt}
\begin{tabularx}{\linewidth}{@{}L{3.2cm} Y@{}}
\toprule
\rowcolor{lgClaro}
\multicolumn{2}{@{}l@{}}{\sffamily\bfseries\color{lgPrimario} #1 \quad #2} \\
\midrule
\sffamily\bfseries\small Prioridad   & #3 \\
\sffamily\bfseries\small Descripción & #4 \\
\sffamily\bfseries\small Criterio de aceptación & #5 \\
\sffamily\bfseries\small Origen      & #6 \\
\bottomrule
\end{tabularx}
\par\vspace{8pt}}

% --- Ficha reutilizable para una historia de usuario -----------------------
% Uso: \historia{HU-001}{rol}{funcionalidad}{beneficio}{criterios}
\newcommand{\historia}[5]{%
\par\vspace{4pt}
\begin{tabularx}{\linewidth}{@{}L{3.2cm} Y@{}}
\toprule
\rowcolor{lgClaro}
\multicolumn{2}{@{}l@{}}{\sffamily\bfseries\color{lgPrimario} #1} \\
\midrule
\sffamily\bfseries\small Como        & #2 \\
\sffamily\bfseries\small Quiero      & #3 \\
\sffamily\bfseries\small Para        & #4 \\
\sffamily\bfseries\small Criterios de aceptación & #5 \\
\bottomrule
\end{tabularx}
\par\vspace{8pt}}

\begin{document}
\lgportada

\lgcontrolversiones{
0.1 & \campo{dd-mm-aaaa} & \lgDirector & Plantilla inicial conforme a ISO/IEC/IEEE 29148. \\
}

\begin{porcompletar}[Cómo usar esta plantilla]
Los apartados reproducen la estructura recomendada por ISO/IEC/IEEE 29148, que
Legasoft adopta como referencia (véase \texttt{LGS-INST-001}, sección 8). Los
campos en ámbar se sustituyen con la información obtenida en el levantamiento.
Cada requisito debe ser verificable: si no se puede escribir su criterio de
aceptación, el requisito todavía no está bien definido.
\end{porcompletar}

\tableofcontents
\clearpage

% =============================================================================
\section{Introducción}

\subsection{Propósito}

Este documento especifica los requisitos funcionales y no funcionales del sistema
\campo{nombre del sistema}, desarrollado por Legasoft para \campo{organización
cliente}. Está dirigido al equipo del proyecto, a la parte interesada y a
cualquier persona que deba verificar la correspondencia entre las necesidades
identificadas y la solución construida.

\subsection{Alcance del producto}

\begin{porcompletar}
\campo{Nombre del sistema, qué hace y qué no hace, y beneficios esperados para la
organización cliente.}
\end{porcompletar}

\subsection{Definiciones, acrónimos y abreviaturas}

\begin{center}
\begin{tabularx}{\linewidth}{@{}L{3.4cm} Y@{}}
\toprule
\sffamily\bfseries\color{lgPrimario}Término &
\sffamily\bfseries\color{lgPrimario}Definición \\
\midrule
SRS      & \emph{Software Requirements Specification}: especificación de requisitos de software. \\
RF       & Requisito funcional. \\
RNF      & Requisito no funcional. \\
HU       & Historia de usuario. \\
API      & \emph{Application Programming Interface}: interfaz de programación de aplicaciones. \\
\campo{término} & \campo{definición propia del dominio del cliente} \\
\bottomrule
\end{tabularx}
\captionof{table}{Glosario del documento.}
\end{center}

\subsection{Referencias}

\begin{itemize}
  \item ISO/IEC/IEEE 29148 — Ingeniería de requisitos.
  \item ISO/IEC 25010 — Modelo de calidad del producto de software.
  \item \texttt{LGS-INST-001} — Constitución y estructura organizacional de Legasoft.
  \item \texttt{LGS-INST-002} — Acta de constitución del proyecto.
  \item \campo{Actas de reunión y entrevistas que originan los requisitos.}
\end{itemize}

\subsection{Visión general del documento}

La sección 2 describe el contexto y las restricciones generales del producto. La
sección 3 detalla los requisitos específicos. La sección 4 presenta las historias
de usuario con sus criterios de aceptación y la sección 5 la matriz de
trazabilidad que conecta necesidad, requisito y verificación.

% =============================================================================
\section{Descripción general}

\subsection{Perspectiva del producto}

\begin{porcompletar}
\campo{¿El sistema es autónomo, reemplaza a otro existente o se integra con
plataformas del cliente? Indicar los sistemas con los que intercambia
información.}
\end{porcompletar}

\subsection{Funciones principales del producto}

\begin{porcompletar}
\campo{Resumen de las capacidades del sistema en un nivel alto, antes del detalle
de la sección 3.}
\end{porcompletar}

\subsection{Características de los usuarios}

\begin{center}
\begin{tabularx}{\linewidth}{@{}L{3.6cm} L{3.6cm} Y@{}}
\toprule
\sffamily\bfseries\color{lgPrimario}Tipo de usuario &
\sffamily\bfseries\color{lgPrimario}Perfil técnico &
\sffamily\bfseries\color{lgPrimario}Uso previsto del sistema \\
\midrule
\campo{rol} & \campo{nivel} & \campo{tareas que realiza} \\
\campo{rol} & \campo{nivel} & \campo{tareas que realiza} \\
\bottomrule
\end{tabularx}
\captionof{table}{Perfiles de usuario del sistema.}
\end{center}

\subsection{Restricciones generales}

\begin{porcompletar}
\campo{Tecnologías obligatorias, normativa aplicable, límites de infraestructura,
políticas de seguridad o capacidad de adopción del cliente.}
\end{porcompletar}

\subsection{Suposiciones y dependencias}

\begin{porcompletar}
\campo{Condiciones asumidas y dependencias externas cuya variación obligaría a
revisar los requisitos.}
\end{porcompletar}

% =============================================================================
\section{Requisitos específicos}

\subsection{Requisitos funcionales}

Cada requisito se identifica de forma única, se prioriza y se acompaña de un
criterio de aceptación verificable. El campo \emph{Origen} mantiene la
trazabilidad con la fuente que lo motivó.

\reqfun{RF-001}{\campo{nombre corto del requisito}}
  {\campo{Alta / Media / Baja}}
  {El sistema deberá \campo{comportamiento observable, en una sola frase}.}
  {\campo{Condición comprobable que permite dar el requisito por cumplido.}}
  {\campo{Entrevista, minuta o documento de origen.}}

\reqfun{RF-002}{\campo{nombre corto del requisito}}
  {\campo{Alta / Media / Baja}}
  {El sistema deberá \campo{comportamiento observable}.}
  {\campo{Criterio de aceptación.}}
  {\campo{Origen.}}

\begin{nota}[Redacción de requisitos]
Un requisito bien formulado es unívoco, verificable, necesario y libre de
detalles de implementación. Evite «el sistema debería», «de forma amigable» o
«rápido»: sustitúyalos por condiciones medibles.
\end{nota}

\subsection{Requisitos no funcionales}

Los atributos se organizan según las características de calidad de ISO/IEC 25010,
norma que Legasoft adopta como referencia para definir y evaluar la calidad del
producto.

\begin{center}
\begin{tabularx}{\linewidth}{@{}C{1.9cm} L{3.4cm} Y@{}}
\toprule
\sffamily\bfseries\color{lgPrimario}ID &
\sffamily\bfseries\color{lgPrimario}Característica (25010) &
\sffamily\bfseries\color{lgPrimario}Requisito y métrica de verificación \\
\midrule
RNF-001 & Eficiencia de desempeño & \campo{p. ej. el 95\% de las consultas responde en menos de 2 s con N usuarios concurrentes.} \\
RNF-002 & Seguridad               & \campo{control de acceso, cifrado, registro de auditoría.} \\
RNF-003 & Usabilidad              & \campo{criterio observable de facilidad de uso o accesibilidad.} \\
RNF-004 & Fiabilidad              & \campo{disponibilidad esperada y comportamiento ante fallos.} \\
RNF-005 & Mantenibilidad          & \campo{modularidad, cobertura de pruebas, documentación exigida.} \\
RNF-006 & Compatibilidad          & \campo{navegadores, dispositivos o sistemas con los que debe operar.} \\
\bottomrule
\end{tabularx}
\captionof{table}{Requisitos no funcionales por característica de calidad.}
\end{center}

\subsection{Requisitos de interfaz}

\subsubsection{Interfaces de usuario}
\begin{porcompletar}
\campo{Pantallas principales, flujos de navegación y requisitos de accesibilidad.}
\end{porcompletar}

\subsubsection{Interfaces con otros sistemas}
\begin{porcompletar}
\campo{APIs consumidas o expuestas, formatos de intercambio y protocolos. El
detalle de los contratos se documenta en el documento de arquitectura
(\texttt{LGS-ARQ-001}).}
\end{porcompletar}

\subsubsection{Interfaces de datos}
\begin{porcompletar}
\campo{Entidades principales, volúmenes estimados y requisitos de integridad,
disponibilidad y tratamiento responsable de la información.}
\end{porcompletar}

% =============================================================================
\section{Historias de usuario y criterios de aceptación}

Las historias de usuario expresan la necesidad desde la perspectiva de quien usa
el sistema y se acuerdan con el cliente. Los criterios de aceptación son la base
sobre la que el Especialista en QA construye las pruebas correspondientes.

\historia{HU-001}
  {\campo{rol de usuario}}
  {\campo{acción que desea realizar}}
  {\campo{beneficio que obtiene}}
  {\campo{Dado \ldots\ cuando \ldots\ entonces \ldots}}

\historia{HU-002}
  {\campo{rol de usuario}}
  {\campo{acción que desea realizar}}
  {\campo{beneficio que obtiene}}
  {\campo{Dado \ldots\ cuando \ldots\ entonces \ldots}}

% =============================================================================
\section{Matriz de trazabilidad}

La matriz conecta la necesidad identificada con el requisito que la recoge, el
elemento de diseño que la implementa y la prueba que la verifica. Es el
instrumento que sostiene el criterio institucional de trazabilidad entre
necesidad, solución y calidad comprobada.

\begin{center}
\begin{tabularx}{\linewidth}{@{}C{1.9cm} L{3.2cm} Y C{2.2cm}@{}}
\toprule
\sffamily\bfseries\color{lgPrimario}Requisito &
\sffamily\bfseries\color{lgPrimario}Necesidad de origen &
\sffamily\bfseries\color{lgPrimario}Elemento de diseño &
\sffamily\bfseries\color{lgPrimario}Caso de prueba \\
\midrule
RF-001 & \campo{origen} & \campo{componente o módulo} & \campo{CP-001} \\
RF-002 & \campo{origen} & \campo{componente o módulo} & \campo{CP-002} \\
RNF-001 & \campo{origen} & \campo{decisión arquitectónica} & \campo{CP-0xx} \\
\bottomrule
\end{tabularx}
\captionof{table}{Trazabilidad entre necesidades, requisitos, diseño y pruebas.}
\end{center}

% =============================================================================
\section{Aprobación del documento}

\vspace{1.0cm}
\begin{center}
\begin{tabularx}{\linewidth}{@{}Y Y@{}}
\centering\rule{5.4cm}{0.5pt} &
\centering\arraybackslash\rule{5.4cm}{0.5pt} \\[-2pt]
\centering\sffamily\small\lgDirector &
\centering\arraybackslash\sffamily\small\campo{representante del cliente} \\[-4pt]
\centering\scriptsize\color{lgGris}Elaboró — \lgDirectorRol &
\centering\arraybackslash\scriptsize\color{lgGris}Aprobó — organización cliente \\
\end{tabularx}
\end{center}

\end{document}
